% The C++ chapter of the WFDB-x user's manual. It follows the sections of the Python
% chapter (chapters/python.tex), in the same order. Each \apientry of the section "The
% high-level API" is a function declared at namespace scope in
% cpp/include/wfdbx/record.hpp; its signature listing is copied from that declaration,
% and its params name every parameter of the declaration, in order. The examples of style
% "example" were compiled and run (each as the body of main(), after the lines given in
% the conventions of the section, in a copy of tests/data); check_manual.py does not run
% them.
\chapter{C++}
\label{chap:cpp}
\chapterlanguage{cpp}

The C++ implementation of WFDB-x is the library \code{libwfdbx}, with its headers in
\file{cpp/include/wfdbx} and everything in namespace \code{wfdbx}, and the 33
command-line tools, built as separate programs. It is a complete C++17 implementation of
the WFDB Software Package: it needs only the C++17 standard library, and never calls the
C library.

It offers WFDB-x in three ways:
\begin{itemize}
\item \textbf{The high-level API} (\cref{sec:cpp-api}): functions, declared in
  \file{<wfdbx/record.hpp>}, that read and write whole records, headers, annotation files
  and header metadata as \code{std::vector}s and structures, in one call, with options
  given as fields of a structure, and that throw \code{wfdbx::Error} on errors. They read
  signal files in large blocks, and are the fastest way to read a record. They are an
  addition of WFDB-x: the C package has no equivalent. They are the recommended interface
  for new code and analysis.
\item \textbf{The WFDB library functions} under their C names (\code{isigopen},
  \code{getvec}, \code{getann}, \code{putann}, \ldots), declared in
  \file{<wfdbx/wfdb.hpp>}, with the behavior of the C library (\cref{sec:cpp-library}).
  They are documented once, for all the languages, in \cref{chap:library}.
\item \textbf{The command-line tools} (\code{rdsamp}, \code{rdann}, \code{wrsamp},
  \code{xform} and 29 others), with the options, output and exit status of the C tools,
  one program per tool, run by its name (\cref{sec:cpp-tools}). They are documented in
  \cref{chap:tools}.
\end{itemize}
The high-level functions return exactly the samples and annotations that the library
functions (and the C library) give.

% ==========================================================================================
\section{Installation and requirements}
\label{sec:cpp-install}

\subsection{Requirements}

WFDB-x for C++ builds with CMake 3.14 or later and a C++17 compiler whose standard
library has \code{std::filesystem}: GCC 9, Clang 9 or MSVC 2019 and later (macOS 10.15 or
later), on Linux, macOS and Windows. It has no other dependency: no WFDB library and no
libFLAC (the FLAC formats 508, 516 and 524 are read and written by WFDB-x's own codec).
Network access (\cref{sec:network}) runs the program \cmd{curl}, which must then be
installed. The unit tests of \code{check\_record} also run \cmd{python3}
(\cref{sec:cpp-testing}).

\subsection{Building and installing}

From the root of a copy of the repository, these commands build the library, the tools
and the unit tests (a Release build), and install them:

\begin{lstlisting}[style=shell]
cmake -S cpp -B cpp/build                        # a Release build
cmake --build cpp/build -j8
cmake --install cpp/build --prefix /usr/local    # optional
\end{lstlisting}

The library is \file{cpp/build/libwfdbx.a}, and the tools are in \file{cpp/build/bin/}.
The install step copies the library, the headers (\file{include/wfdbx/}), the tools and a
CMake package file (\file{lib/cmake/wfdbx/wfdbxConfig.cmake}, with the target
\code{wfdbx::wfdbx}). These options of the build change what is built:

\begin{table}[htbp]
\centering
\caption{The options of the CMake build}
\label{tab:cpp-cmake-options}
\begin{tabularx}{\linewidth}{@{}llL@{}}
\toprule
\tabhead Option & \tabhead Default & \tabhead Effect \\
\midrule
\code{WFDBX\_BUILD\_TOOLS} & \code{ON} & build the command-line tools \\
\code{WFDBX\_BUILD\_TESTS} & \code{ON} & build the unit tests \\
\code{BUILD\_SHARED\_LIBS} & \code{OFF} & build a shared library \\
\code{CMAKE\_BUILD\_TYPE} & \code{Release} & the build type, when none is given \\
\bottomrule
\end{tabularx}
\end{table}

The version of the library is that of the CMake project, \code{1.2.0}. The tools, like
the C tools, have no \opt{-{}-version} option.

\subsection{Using the library from CMake}

An installed library is found with \code{find\_package}, which defines the imported
target \code{wfdbx::wfdbx} (the library and its include directory):

\begin{lstlisting}
find_package(wfdbx REQUIRED)
target_link_libraries(myprogram PRIVATE wfdbx::wfdbx)
\end{lstlisting}

A copy of the directory \file{cpp} (for example in a subdirectory of a project) can also
be built as part of a project with \code{add\_subdirectory}, which defines the target
\code{wfdbx} and its alias \code{wfdbx::wfdbx}. The options of
\cref{tab:cpp-cmake-options} then select what is built; for the library alone:

\begin{lstlisting}
set(WFDBX_BUILD_TOOLS OFF CACHE BOOL "Build the command-line tools" FORCE)
set(WFDBX_BUILD_TESTS OFF CACHE BOOL "Build the unit tests" FORCE)
add_subdirectory(path/to/wfdb-x/cpp ${CMAKE_CURRENT_BINARY_DIR}/wfdbx EXCLUDE_FROM_ALL)
target_link_libraries(myprogram PRIVATE wfdbx::wfdbx)
\end{lstlisting}

Without CMake, a program is compiled against the headers and linked with the library,
installed (\cmd{-lwfdbx}) or in the build directory:

\begin{lstlisting}[style=shell]
c++ -std=c++17 myprogram.cpp -lwfdbx                                  # installed
c++ -std=c++17 -Icpp/include myprogram.cpp cpp/build/libwfdbx.a -o myprogram
\end{lstlisting}

\subsection{The sample record}
\label{sec:cpp-resources}

The repository includes the sample record \code{100s}, the sample record of the WFDB
Software Package: the first minute of record 100 of the MIT-BIH Arrhythmia Database, two
ECG signals (\code{MLII} and \code{V5}) sampled at 360\,Hz, with the reference beat
annotations (annotator \code{atr}). Its files (\file{100s.hea}, \file{100s.dat},
\file{100s.atr}) are in the directory \file{tests/data} of a copy of the repository, with
the other test records of WFDB-x. The C++ library does not install them: programs name
that directory as their input directory (\code{wfdbx\_set\_input\_dir("tests/data")}).
Work that uses the MIT-BIH Arrhythmia Database cites it as described on PhysioNet.

% ==========================================================================================
\section{Quick start}
\label{sec:cpp-quickstart}

This program reads the header, ten seconds of the signals and the annotations of the
sample record \code{100s}, from the directory \file{tests/data} of the repository (the
input directory of each call), and
computes the mean heart rate. It then writes the first ten seconds as a new record,
\code{first10}, in the current directory, with a metadata entry in its header.

\begin{lstlisting}[language=C++,style=example]
#include <cstdio>
#include <vector>
#include <wfdbx/ecgmap.hpp>
#include <wfdbx/record.hpp>

int main() {
    wfdbx::InputOptions sample;
    sample.input_dir = "tests/data";             // the sample record 100s

    // The header only: no samples are read.
    wfdbx::Header hdr = wfdbx::read_header("100s", false, sample);
    std::printf("%d %g %lld\n", hdr.nsig, hdr.fs, hdr.nsamples);   // 2 360 21600
    std::printf("%s %s\n", hdr.siginfo[0].desc->c_str(),
                hdr.siginfo[1].desc->c_str());                     // MLII V5

    // Seconds 10 to 20 of both signals, frame by frame, in ADC units.
    wfdbx::ReadOptions opt;
    opt.input_dir = "tests/data";
    opt.from = "0:10";
    opt.to = "0:20";
    wfdbx::Record rec = wfdbx::read_record("100s", opt);
    std::printf("%zu %d\n", rec.nsamples(), rec.nsig);             // 3600 2
    std::printf("%d %g\n", rec.at(0, 0), rec.physical(0, 0));      // 946 -0.39

    // The beat annotations, and the mean heart rate.
    wfdbx::Annotations ann = wfdbx::read_annotations("100s", "atr", {}, {},
                                                     std::nullopt, false, sample);
    std::vector<long long> beats;
    for (size_t i = 0; i < ann.size(); i++)
        if (wfdbx::isqrs(ann.codes[i])) beats.push_back(ann.samples[i]);
    double rr = (beats.back() - beats.front()) / ann.fs / (beats.size() - 1);
    std::printf("%zu beats, %.1f beats per minute\n", beats.size(), 60 / rr);
    // 74 beats, 73.9 beats per minute

    // Write the first 10 seconds as the record "first10" (first10.hea and first10.dat,
    // in the current directory), with the signal specifications of 100s.
    opt.from = 0;
    opt.to = "0:10";
    wfdbx::Record x = wfdbx::read_record("100s", opt);
    wfdbx::WriteOptions w;
    w.gains = {200};
    w.baselines = {1024};
    w.adczeros = {1024};
    w.adcres = {11};
    w.units = {"mV"};
    w.names = x.names();
    w.metadata = {{"Source", "MIT-BIH record 100"}};
    wfdbx::write_record("first10", x.samples, x.nsig, x.fs, w);

    // Header metadata: "# Key: value" lines of the header.
    wfdbx::set_metadata("first10", "Age", "69");
    for (const auto &e : wfdbx::read_metadata("first10"))
        std::printf("%s = %s\n", e.first.c_str(), e.second.c_str());
    // Source = MIT-BIH record 100
    // Age = 69
    std::printf("%s\n", wfdbx::check_record("first10").to_string().c_str());
    // no problems found
}
\end{lstlisting}

Compiled and run from the root of the repository:

\begin{lstlisting}[style=shell]
c++ -std=c++17 -Icpp/include quick.cpp cpp/build/libwfdbx.a -o quick
./quick
\end{lstlisting}

The header \file{100s.hea} gives no units: the program writes \code{"mV"} explicitly,
since WFDB-x never assumes units, a gain, an ADC resolution or a sampling frequency that
a record does not give (\cref{chap:behavior}).

% ==========================================================================================
\section{The high-level API}
\label{sec:cpp-api}

The high-level API is made of the functions of this section, declared in
\file{<wfdbx/record.hpp>} in namespace \code{wfdbx}, of the structures of their options
(\cref{sec:cpp-options}) and of their results (\cref{sec:cpp-results}), and of the
exceptions \code{Error}, \code{ArgumentError} and \code{FileError}
(\cref{sec:cpp-warnings}). The comments of the header document each of them.

\subsection{Conventions}
\label{sec:cpp-conventions}

\paragraph{Record names and directories.} A record is named as in the other parts of
WFDB-x (\cref{chap:records,chap:files}), without the \file{.hea} extension (the readers
also accept a trailing \code{".hea"}; \code{write\_record} does not). A record name is
one of: a full path, absolute (\code{"/data/mitdb/100"}) or in the home directory
(\code{"\textasciitilde/Documents/100s"}), used as given; a relative path
(\code{"100s"}, \code{"data/100s"}), read from the input directory if one is set,
otherwise from the current directory, and written to the output directory if one is set,
otherwise to the current directory (missing directories are created); or a URL
(\code{"https://physionet.org/files/mitdb/1.0.0/100"}), read only if network access is
on, and which cannot be written. The input and output directories are set with
\code{wfdbx\_set\_input\_dir} and \code{wfdbx\_set\_output\_dir}
(\cref{sec:cpp-settings}) or with the environment variables \code{WFDBX\_INPUT\_DIR} and
\code{WFDBX\_OUTPUT\_DIR}. A writer given a name with a path creates the files in that
directory, and the header names the record and its signal file without the path
(\code{rec}, \code{rec.dat}), so that the record opens under any path that leads to it;
in legacy mode, as in the original, the header gives the name with its path. The readers
take the fields \code{input\_dir} and \code{network} of an \code{InputOptions} (or of a
structure derived from it), and the writers a \code{dir}: they apply to that one call,
and the previous settings are restored when it returns, also when it throws. The
\code{input\_dir} option is optional, and applies to relative names only, not to full
paths or URLs. With \code{dir}, the record name must be a plain name (\code{"rec"}): the
files are created in \code{dir}, and the record opens from that directory under its plain
name. The signal files that a header names are read from the header's directory.

\paragraph{Options.} The options of a function are the fields of a structure
(\code{ReadOptions}, \code{InputOptions}, \code{WriteOptions}, \code{MetadataOptions}),
given as its last argument, with a default value \code{\{\}}: a field that is not set
keeps its default. An option that may be absent is a \code{std::optional}.

\paragraph{Times.} The fields \code{from} and \code{to} of \code{ReadOptions}, and the
arguments \code{from} and \code{to} of \code{read\_annotations}, are \code{TimeSpec}
values: a sample number or a WFDB time string, as \code{strtim} reads it (\code{"0:10"},
10 seconds; \code{"1:02:03.5"}; \code{"s3600"}, sample 3600; \code{"e"}, the end of the
record; or a time of day in brackets, \code{"[14:00:00]"}). A \code{TimeSpec} converts
from integers and strings, so \code{opt.from = 3600;} and \code{opt.from = "0:10";} both
work; a default-constructed \code{TimeSpec} is empty (no time is given). A time in seconds
needs the sampling frequency; when the record does not give it, such a time throws
\code{Error}.

\paragraph{Samples.} Samples are \code{WFDB\_Sample} values (\code{int}), in ADC units,
stored frame by frame in one \code{std::vector}: \code{samples[i * nsig + s]} is sample
\code{i} of signal \code{s}. Invalid samples are \code{WFDB\_INVALID\_SAMPLE}
($-32768$), and NaN in physical units.

\paragraph{Per-signal and per-annotation values.} The per-signal fields of
\code{WriteOptions} (\code{gains}, \code{units}, \ldots) hold one value for all the
signals, or one value per signal; the per-annotation fields of \code{Annotations}
(\code{symbols}, \code{chans}, \ldots) are empty or hold one value per annotation. Any
other length throws \code{Error} before anything is written.

\paragraph{Signal parameters.} No sampling frequency, gain, ADC resolution or units is
ever assumed. Where one is needed and the record does not give it, it must be given;
otherwise the function throws \code{Error}. A record whose header gives no sampling
frequency has \code{fs = 0}.

\paragraph{Text.} Strings are \code{std::string}, handled as bytes: every byte sequence
round-trips unchanged, and UTF-8 text is stored as its bytes.

\paragraph{Legacy mode.} With \code{legacy} set to \code{true}, a function behaves as in
the original WFDB package, with its assumed signal parameters and its search for files
(\cref{chap:behavior}); the library's legacy mode is on for the duration of the call
only. With \code{false} (the default), the current setting of \code{wfdbx\_set\_legacy}
applies, which is off unless it was turned on.

\paragraph{State.} The functions that read and write records and annotation files call
\code{wfdbquit()} before and after their work, also when they throw, which closes
anything opened with the WFDB library functions. The library's settings (legacy mode,
getvec mode, time resolution of output annotators, input and output directories) are the
same after a call as before it. Opening an EDF record sets the time resolution of output
annotators (\code{getafreq()}) to its sampling frequency when none is set, and adds
\code{WFDB\_HIGHRES} to the getvec mode (\code{getgvmode()}), as in C, and
\code{wfdbquit()} resets neither: the functions that open a record restore both when they
return, also when they throw, so that a later \code{write\_annotations} does not record
that frequency; the C-style functions keep the behavior of C. The metadata functions do
not call \code{wfdbquit()}: they open only the header, and leave the library's open files
as they are. The functions are not thread-safe. Results are returned by value.

\paragraph{The examples} of this section are the body of a function \code{main}, compiled
after these lines, and run in a copy of the test records of \file{tests/data} (among them
\code{100s}), which is the current directory and holds the records that the examples
write:
\begin{lstlisting}[language=C++]
#include <algorithm>
#include <cmath>
#include <cstdio>
#include <filesystem>
#include <random>
#include <vector>
#include <wfdbx/ecgcodes.hpp>
#include <wfdbx/record.hpp>
using namespace wfdbx;
\end{lstlisting}

% ------------------------------------------------------------------------------------------
\subsection{Reading records}
\label{sec:cpp-reading}

\apientry{read\_record}{Read the signals of a record}
\begin{lstlisting}[language=C++,style=signature]
Record read_record(const std::string &record, const ReadOptions &opt = {});
\end{lstlisting}
Reads the signals of a record from \code{opt.from} up to \code{opt.to} (exclusive), in
ADC units, frame by frame, together with the description of the record. The samples are
those that \code{getvec} returns (\cref{chap:library}), read in large blocks with
\code{getvecs}. If the record ends before \code{opt.to}, the samples up to its end are
returned. When the read reaches the end of the record, the samples are compared with the
checksums of the header (\cref{sec:cpp-warnings}). The getvec mode (\code{getgvmode()})
and the time resolution of output annotators (\code{getafreq()}) are the same after the
call as before it (opening an EDF record sets both).
\begin{params}
\item[record] Record name, without the \file{.hea} extension (a trailing \code{".hea"} is
  also accepted): a full path, absolute (\code{"/data/mitdb/100"}) or in the home
  directory (\code{"\textasciitilde/Documents/100s"}), used as given; a relative path
  (\code{"100s"}, \code{"data/100s"}), read from the input directory if one is set,
  otherwise from the current directory; or a URL
  (\code{"https://physionet.org/files/mitdb/1.0.0/100"}), read only if network access is
  on.
\item[opt] The options (\cref{tab:cpp-readoptions}): the signals to read
  (\code{channels}), the interval (\code{from}, \code{to}), \code{highres}, \code{pad},
  \code{checksum\_warning}, \code{legacy}, and the input directory, under which relative
  record names are read, and network access of the call (\code{input\_dir},
  \code{network}). Default \code{\{\}}: every signal of the
  whole record, with the current settings.
\end{params}
\returns{A \code{Record} (\cref{tab:cpp-record}), with the samples in ADC units.}
\raises{\code{Error} if the record cannot be read (``can't read record REC'') or has no
signals; if a requested signal does not exist (``record REC has no signal C''); if
\code{opt.from} is negative (``read\_record: from must not be negative'') or cannot be
reached (``can't seek to sample N of record REC''); if a time string needs the sampling
frequency and the header gives none (``sampling frequency of record REC is not
specified''); or if a signal file is shorter than its header says, or cannot be read or
decoded (``can't read signal data of record REC (getvec status S)''). Data that do not
match the checksums of the header are returned, with a warning (see
\code{checksum\_warning}).}
\begin{lstlisting}[language=C++,style=example]
// Seconds 10 to 20 of lead MLII, in physical units.
ReadOptions opt;
opt.channels = {"MLII"};
opt.from = "0:10";
opt.to = "0:20";
Record rec = read_record("100s", opt);
std::vector<double> mlii = rec.physical_signal(0);
std::printf("%zu %g %lld %s\n", mlii.size(), rec.fs, rec.sampfrom,
            rec.names()[0].c_str());                   // 3600 360 3600 MLII

// The first 3 samples of both signals, in ADC units, V5 first.
ReadOptions first;
first.channels = {"1", "0"};
first.to = 3;
rec = read_record("100s", first);
for (size_t i = 0; i < rec.nsamples(); i++)
    std::printf("%d %d\n", rec.at(i, 0), rec.at(i, 1));   // 1011 995, 3 times
std::printf("%g %d\n", rec.siginfo[0].gain, rec.siginfo[0].baseline);   // 200 1024

// The whole record: the checksums are verified.
rec = read_record("100s");
std::printf("%zu %d\n", rec.nsamples(), rec.checksum_ok);   // 21600 1
\end{lstlisting}
A multi-frequency record (here the test record \code{mfreq}, whose signals have 1, 2
and 4 samples per frame at a frame frequency of 100\,Hz) is read at its frame frequency,
or at its highest sampling frequency:
\begin{lstlisting}[language=C++,style=example]
Record rec = read_record("mfreq");
std::printf("%g %zu\n", rec.fs, rec.nsamples());       // 100 600
ReadOptions opt;
opt.highres = true;
rec = read_record("mfreq", opt);
std::printf("%g %zu\n", rec.fs, rec.nsamples());       // 400 2400
for (const WFDB_Siginfo &s : rec.siginfo) std::printf("%d ", s.spf);   // 1 2 4
std::printf("\n");
\end{lstlisting}

\apientry{read\_header}{Read the header of a record}
\begin{lstlisting}[language=C++,style=signature]
Header read_header(const std::string &record, bool legacy = false,
                   const InputOptions &input = {});
\end{lstlisting}
Reads the header of a record, without opening its signal files. The result describes the
record and each of its signals, and holds the info strings and the metadata entries of the
header (from the \file{.hea} file only, as \code{read\_metadata} reads them, with its
warning for a key given more than once). For an EDF record, which has no \file{.hea}
file, and for a header read from the standard input (\code{"-"}), the metadata is empty,
without an error. For a multi-segment record, the master header is read, and
\code{nsamples} is the length of the whole record. The time resolution of output
annotators (\code{getafreq()}) and the getvec mode (\code{getgvmode()}) are the same after
the call as before it (opening an EDF record sets both).
\begin{params}
\item[record] Record name, without the \file{.hea} extension (a trailing \code{".hea"} is
  also accepted): a full path, absolute (\code{"/data/mitdb/100"}) or in the home
  directory (\code{"\textasciitilde/Documents/100s"}), used as given; a relative path
  (\code{"100s"}, \code{"data/100s"}), read from the input directory if one is set,
  otherwise from the current directory; or a URL
  (\code{"https://physionet.org/files/mitdb/1.0.0/100"}), read only if network access is
  on.
\item[legacy] \code{true} to read the header as the original WFDB package does. Default
  \code{false}.
\item[input] The input directory for this call only (optional), under which relative
  record names are read (it does not apply to full paths or URLs), and network access
  (\cref{tab:cpp-inputoptions}). Default \code{\{\}}: the current settings.
\end{params}
\returns{A \code{Header} (\cref{tab:cpp-header}).}
\raises{\code{Error} if the header cannot be read (``can't read header for record
REC'').}
\begin{lstlisting}[language=C++,style=example]
Header hdr = read_header("100s");
std::printf("%s %g %d %lld\n", hdr.name.c_str(), hdr.fs, hdr.nsig, hdr.nsamples);
// 100s 360 2 21600
for (const WFDB_Siginfo &s : hdr.siginfo)
    std::printf("%s %s %d %g %d %s\n", s.desc->c_str(), s.fname->c_str(), s.fmt,
                s.gain, s.adcres, s.units ? s.units->c_str() : "(none)");
// MLII 100s.dat 212 200 11 (none)
// V5 100s.dat 212 200 11 (none)
std::printf("[%s]\n", hdr.comments[0].c_str());   // [ 69 M 1085 1629 x1]
std::printf("%d\n", hdr.base_time.has_value());   // 0: the header gives no start time
\end{lstlisting}

\apientry{read\_annotations}{Read an annotation file}
\begin{lstlisting}[language=C++,style=signature]
Annotations read_annotations(const std::string &record,
                             const std::string &annotator,
                             const TimeSpec &from = {},
                             const TimeSpec &to = {},
                             std::optional<WFDB_Frequency> fs = std::nullopt,
                             bool legacy = false,
                             const InputOptions &input = {});
\end{lstlisting}
Reads the annotations of an annotation file, in file order, as parallel vectors. The
sampling frequency, needed to convert times in seconds, is that of the record's header,
or, for a record without a header or whose header gives none, the argument \code{fs}, as
the option \opt{-F} of \code{rdann}. The time resolution of output annotators
(\code{getafreq()}) and the getvec mode (\code{getgvmode()}) are the same after the call
as before it (opening an EDF record sets both).
\begin{params}
\item[record] Record name, without the \file{.hea} extension (a trailing \code{".hea"} is
  also accepted): a full path, absolute (\code{"/data/mitdb/100"}) or in the home
  directory (\code{"\textasciitilde/Documents/100s"}), used as given; a relative path
  (\code{"100s"}, \code{"data/100s"}), read from the input directory if one is set,
  otherwise from the current directory; or a URL
  (\code{"https://physionet.org/files/mitdb/1.0.0/100"}), read only if network access is
  on. Its header, if there is one, gives the sampling
  frequency.
\item[annotator] Annotator name: the suffix of the annotation file. \code{"atr"} reads the
  file \file{RECORD.atr}.
\item[from] The first time to read: a sample number or a time string. Default \code{\{\}}:
  the beginning of the file.
\item[to] The time at which to stop: annotations at or after it are not read. Default
  \code{\{\}}: the end of the file.
\item[fs] The sampling frequency, in Hz ($> 0$), to use when the record has no header or
  its header gives none; it is ignored otherwise. Default \code{std::nullopt}: none.
\item[legacy] \code{true} to behave as the original WFDB package (which assumes 250\,Hz
  for a record without a frequency). Default \code{false}.
\item[input] The input directory for this call only (optional), under which relative
  record names are read (it does not apply to full paths or URLs), and network access
  (\cref{tab:cpp-inputoptions}). Default \code{\{\}}: the current settings.
\end{params}
\returns{An \code{Annotations} (\cref{tab:cpp-annotations}). Its \code{fs} is the
frequency in which the times are counted: the one written in the annotation file, if
there is one, or else the sampling frequency; it is 0 if none of them is known.}
\raises{\code{Error} if \code{fs} is given and is not $> 0$ (``read\_annotations: fs must
be > 0''); if the annotation file cannot be opened (``can't read annotator ANN for record
REC''); if a time string needs the sampling frequency and none is known (``sampling
frequency of record REC is not specified''); or if the file ends early or cannot be read
before \code{to} is reached (``can't read annotation file ANN of record REC (getann status
S)'', where S is the status that \code{getann} returned, such as $-3$ for an unexpected
end of file).}
\begin{lstlisting}[language=C++,style=example]
Annotations ann = read_annotations("100s", "atr", "0:10", "0:30");
std::printf("%zu %g\n", ann.size(), ann.fs);                   // 24 360
std::printf("%lld %s\n", ann.samples[0], ann.symbols[0].c_str());   // 3862 N

// The whole file: the first annotation is a rhythm change, with an aux string.
ann = read_annotations("100s", "atr");
std::printf("%s %d %s\n", ann.symbols[0].c_str(), ann.codes[0] == RHYTHM,
            ann.aux[0]->c_str());                              // + 1 (N
\end{lstlisting}
For a record without a header, \code{fs} gives the frequency in which times in seconds
are converted (the library reports on the standard error that it cannot open the
header):
\begin{lstlisting}[language=C++,style=example]
Annotations beats;
beats.samples = {100, 460, 820};
beats.symbols = {"N", "N", "N"};
write_annotations("beats", "qrs", beats);             // no header
Annotations ann = read_annotations("beats", "qrs", {}, "0:02", 360.0);
std::printf("%zu %g\n", ann.size(), ann.fs);          // 2 360
\end{lstlisting}

% ------------------------------------------------------------------------------------------
\subsection{Writing records}
\label{sec:cpp-writing}

\apientry{write\_record}{Write a header and a signal file}
\begin{lstlisting}[language=C++,style=signature]
void write_record(const std::string &record,
                  const std::vector<WFDB_Sample> &samples, int nsig,
                  WFDB_Frequency fs, const WriteOptions &opt = {});
\end{lstlisting}
Writes samples in ADC units as a new record: the header \file{RECORD.hea} and one signal
file, \file{RECORD.dat}, which holds all the signals (so they share one format). Existing
files of these names are overwritten. The sampling frequency, the gain and ADC resolution
of each signal, and the units of each calibrated signal must be given: there are no
defaults. Every check except those of the writing itself comes before any file is
created. The time resolution of output annotators (\code{getafreq()}) and the getvec
mode (\code{getgvmode()}) are left as they were.
\begin{params}
\item[record] Record name (not empty), without the \file{.hea} extension: a plain name
  (\code{"rec"}) or a relative path (\code{"data/rec"}), created in the output directory
  if one is set, otherwise in the current directory, missing directories included; or a
  full path, absolute or in the home directory (\code{"\textasciitilde/Documents/rec"}),
  used as given. With a path, the files are created in that directory, and the header
  names the record and its signal file without the path (\code{rec}, \code{rec.dat}), so
  that the record opens under any path that leads to it (in legacy mode, as in the
  original, the header gives the name with its path). A URL cannot be written.
\item[samples] The samples, in ADC units, frame by frame: \code{nsig} values per frame,
  or, for a multi-frequency record, the sum of \code{opt.samps\_per\_frame} (the samples
  of signal 0, then those of signal 1, and so on). \code{WFDB\_INVALID\_SAMPLE} marks
  invalid samples. The number of samples must be a multiple of the frame size.
\item[nsig] The number of signals ($> 0$).
\item[fs] The sampling frequency, in Hz: a finite number $> 0$ (for a multi-frequency
  record, the frame frequency).
\item[opt] The options (\cref{tab:cpp-writeoptions}): the format, the signal
  specifications (\code{gains}, \code{adcres} and \code{units} are required), the start
  time, the comments and metadata, the directory \code{dir} (which must exist) in which
  to create the files for this call only, with a plain record name, and \code{legacy}.
\end{params}
\returns{Nothing.}
\raises{\code{Error} if a \code{samps\_per\_frame} value is less than 1
(``write\_record: samps\_per\_frame must hold integers >= 1'', checked first); if the
record name is empty; if \code{nsig} is not $> 0$ (``write\_record: at least one signal
is required''); if a per-signal field of \code{opt} has another length than 0, 1 or
\code{nsig}; if \code{fs} is not finite or not $> 0$; if a gain is not finite
(``write\_record: gains must be finite''); if the gains, the ADC resolutions or the units
of a calibrated signal are missing (``write\_record: gains must be specified'',
``write\_record: ADC resolution(s) must be specified'', ``write\_record: units must be
specified''); if \code{samples.size()} is not a multiple of the frame size; if
\code{opt.base\_time} is not a valid time (``write\_record: incorrect base time 'T'''); if
the files cannot be created; or if they cannot be written completely, for example on a
full disk (``write error on record REC''). With \code{opt.dir}, also if the directory does
not exist (``directory D does not exist'') or the record name contains one (``record name
must not contain a directory when dir is given''). \code{ArgumentError} (a subclass of
\code{Error}) if an entry of \code{opt.metadata} is not valid (``write\_record: metadata
key 'K' is also given in comments'', for example).}
\begin{lstlisting}[language=C++,style=example]
// A ramp of 1000 samples of one signal, in millivolts, at 100 Hz: the record "ramp"
// (ramp.hea and ramp.dat).
std::vector<WFDB_Sample> x(1000);
for (int i = 0; i < 1000; i++) x[i] = i;
WriteOptions w;
w.gains = {200};
w.adcres = {12};
w.units = {"mV"};
w.names = {"ramp"};
w.base_time = "14:00:03 21/06/2024";
w.comments = {"a ramp"};
w.metadata = {{"Generator", "loop"}};
write_record("ramp", x, 1, 100, w);
Record rec = read_record("ramp");
std::printf("%zu %s %g\n", rec.nsamples(), rec.base_time->c_str(),
            rec.physical(10, 0));                      // 1000 [14:00:03 21/06/2024] 0.05
for (const std::string &c : read_header("ramp").comments)
    std::printf("[%s]\n", c.c_str());
// [a ramp]
// [ Generator: loop]
\end{lstlisting}
A multi-frequency record: signal \code{fast} has 2 samples per frame, signal
\code{slow} one; each frame holds the two samples of \code{fast} and then the sample of
\code{slow}:
\begin{lstlisting}[language=C++,style=example]
std::vector<WFDB_Sample> frames;
for (int i = 0; i < 100; i++) {          // 1 second at 100 frames per second
    frames.push_back(2 * i);             // fast: 200 samples per second
    frames.push_back(2 * i + 1);
    frames.push_back(500);               // slow: 100 samples per second
}
WriteOptions w;
w.gains = {200};
w.adcres = {12};
w.units = {"mV"};
w.names = {"fast", "slow"};
w.samps_per_frame = {2, 1};
write_record("mf", frames, 2, 100, w);
ReadOptions opt;
opt.highres = true;
Record rec = read_record("mf", opt);
std::printf("%g %d %d %d %d\n", rec.fs, rec.at(0, 0), rec.at(0, 1), rec.at(1, 0),
            rec.at(1, 1));                         // 200 0 500 1 500
\end{lstlisting}

\apientry{write\_record\_physical}{Write a record from samples in physical units}
\begin{lstlisting}[language=C++,style=signature]
void write_record_physical(const std::string &record,
                           const std::vector<double> &samples, int nsig,
                           WFDB_Frequency fs, const WriteOptions &opt = {});
\end{lstlisting}
As \code{write\_record}, with the samples in physical units. Each sample $v$ is stored as
$\mathrm{round}(v \cdot \mathit{gain}) + \mathit{baseline}$ (halves rounded away from
zero), and NaN is stored as an invalid sample. Every signal must have a gain other than
0.
\begin{params}
\item[record] Record name, as for \code{write\_record}.
\item[samples] The samples, in physical units, frame by frame (as for
  \code{write\_record}).
\item[nsig] The number of signals ($> 0$).
\item[fs] The sampling frequency, in Hz: a finite number $> 0$.
\item[opt] The options, as for \code{write\_record} (\cref{tab:cpp-writeoptions}).
\end{params}
\returns{Nothing.}
\raises{As \code{write\_record}; also \code{Error} if a gain is 0 (``write\_record: signal
N is uncalibrated''), or if a sample is infinite or its ADC value is outside the range of
a 32-bit integer (``write\_record: sample I is out of range''). Nothing is written then.}
\begin{lstlisting}[language=C++,style=example]
// Two signals of random numbers, in millivolts, at 360 Hz: the record "rand".
std::mt19937 gen(1);
std::normal_distribution<double> randn;
std::vector<double> x(3600 * 2);                    // 3600 frames of 2 signals
for (double &v : x) v = randn(gen);
WriteOptions w;
w.fmt = 16;
w.gains = {1000};
w.adcres = {16};
w.units = {"mV"};
w.names = {"A", "B"};
write_record_physical("rand", x, 2, 360, w);
Record rec = read_record("rand");
double err = 0;
for (size_t i = 0; i < x.size(); i++)
    err = std::max(err, std::abs(rec.physical(i / 2, i % 2) - x[i]));
std::printf("%zu %d\n", rec.nsamples(), err <= 0.0005);   // 3600 1: rounded to 1/1000 mV
\end{lstlisting}

\apientry{write\_annotations}{Write an annotation file}
\begin{lstlisting}[language=C++,style=signature]
void write_annotations(const std::string &record, const std::string &annotator,
                       const Annotations &ann, WFDB_Frequency fs = 0,
                       const std::string &dir = {},
                       std::optional<bool> sort = std::nullopt);
\end{lstlisting}
Writes an annotation file, \file{RECORD.ANNOTATOR}, in the MIT format; an existing file of
that name is overwritten. The record need not have a header. The type of each annotation
is its \code{symbols} entry (\code{"N"}, \code{"V"}, \ldots; an unknown mnemonic gives
\code{NOTQRS}), or its \code{codes} entry (\cref{chap:formats}). Both may be given (as
\code{read\_annotations} returns them) if they agree: the code of each non-empty symbol
(\code{strann}) is the code given for that annotation; an empty symbol leaves the code.
\code{subtypes}, \code{chans} and \code{nums} are stored in 8 bits, as by \code{putann}
(300 is stored as 44, as in C), and \code{aux} texts are cut to 255 bytes. If the
annotations are not in time order, the file is sorted when it is written. The arguments
are checked before the file is created. The time resolution of output annotators
(\code{getafreq()}) and the getvec mode (\code{getgvmode()}) are left as they were:
\code{fs} applies to this file only.
\begin{params}
\item[record] Record name, without the \file{.hea} extension (a trailing \code{".hea"} is
  also accepted): a plain name or a relative path (\code{"data/100s"}) is created in the
  output directory if one is set, otherwise in the current directory; a full path,
  absolute or in the home directory (\code{"\textasciitilde/Documents/100s"}), is used as
  given. The file is \file{RECORD.ANNOTATOR} at that place (\file{data/100s.qrs}). A URL
  cannot be written.
\item[annotator] Annotator name: the suffix of the file, made of letters, digits,
  \code{\_}, \code{-} and \code{\textasciitilde}.
\item[ann] The annotations (\cref{tab:cpp-annotations}): \code{samples}, and
  \code{symbols} or \code{codes}; \code{subtypes}, \code{chans}, \code{nums} and
  \code{aux} are optional. Each of these vectors is empty (\code{subtypes}, \code{chans}
  and \code{nums} are then 0, and there is no aux) or has one entry per annotation.
  \code{ann.record}, \code{ann.annotator} and \code{ann.fs} are not used.
\item[fs] 0 (none, the default) or a finite frequency $> 0$, in which the times are
  counted, recorded in the file.
\item[dir] If not empty, the directory (which must exist) in which to create the file, for
  this call only; \code{record} must then be a plain name, without a directory: the file
  is \file{DIR/RECORD.ANNOTATOR}. Default \code{\{\}}: as \code{record} says.
\item[sort] \code{true} to always sort the file, \code{false} to keep the given order
  (without the library's advice to run \code{sortann}), for this call only. Default
  \code{std::nullopt}: the environment variable \code{WFDBANNSORT} decides (sorting is on
  unless it is \code{0}).
\end{params}
\returns{Nothing.}
\raises{\code{Error} if a per-annotation field has another length than 0 or the number of
annotations (``write\_annotations: symbols must have one entry per annotation (N), or
none'', for example); if \code{symbols} and \code{codes} are both given and disagree
(``write\_annotations: symbols and codes disagree at annotation K'', K counted from 0);
if \code{fs} is not 0 and not a finite frequency $> 0$; if the file cannot be created
(``can't create annotator ANN for record REC''); if it cannot be written completely
(``write error on annotator ANN for record REC''); or, with \code{dir}, as
\code{write\_record} does.}
\begin{lstlisting}[language=C++,style=example]
// Three beats and a rhythm change, by symbol, with an aux string for the rhythm.
Annotations a;
a.samples = {18, 77, 370, 662};
a.symbols = {"+", "N", "V", "N"};
a.aux = {std::string("(N"), std::nullopt, std::nullopt, std::nullopt};
write_annotations("100s", "test", a);
Annotations b = read_annotations("100s", "test");
std::printf("%s %s %s\n", b.symbols[0].c_str(), b.symbols[2].c_str(),
            b.aux[0]->c_str());                          // + V (N

// A copy of the reference annotations, as read_annotations returns them.
Annotations atr = read_annotations("100s", "atr");
write_annotations("100s", "copy", atr);
Annotations c = read_annotations("100s", "copy");
std::printf("%d\n", c.samples == atr.samples && c.aux == atr.aux);   // 1

// Annotations out of time order are sorted when the file is written (the library
// reports it on the standard error).
Annotations u;
u.samples = {500, 100};
u.codes = {NORMAL, NORMAL};
write_annotations("100s", "unsorted", u);
Annotations s = read_annotations("100s", "unsorted");
std::printf("%lld %lld\n", s.samples[0], s.samples[1]);    // 100 500
\end{lstlisting}

% ------------------------------------------------------------------------------------------
\subsection{Checking records}
\label{sec:cpp-checking}

\apientry{check\_record}{Check a record and its annotation files}
\begin{lstlisting}[language=C++,style=signature]
CheckReport check_record(const std::string &record,
                         const std::vector<std::string> &annotators = {},
                         const InputOptions &input = {});
\end{lstlisting}
Checks the header of a record, its signal specifications, every sample of its signal
files, each segment of a multi-segment record, and the annotation files named, and
reports each problem found. A problem in the record never throws: it is reported, with a
severity (\code{"error"}: the record is damaged or does not follow the WFDB format, and
programs will read wrong data or fail; \code{"warning"}: the record can be read, but
something is missing or unusual, and may give wrong results; \code{"info"}: a fact worth
knowing, such as invalid samples). The problems, their codes and their messages are the
same in the seven languages (\cref{tab:python-issue-codes}). The checks run in normal
mode, even if legacy mode is on, and the library prints no messages during the call (the
previous settings are restored afterwards, the time resolution of output annotators and
the getvec mode among them). They are similar to those of Benjamin Moody's
\code{wfdbcheck}.
\begin{params}
\item[record] Record name, without the \file{.hea} extension (a trailing \code{".hea"} is
  also accepted): a full path, absolute (\code{"/data/mitdb/100"}) or in the home
  directory (\code{"\textasciitilde/Documents/100s"}), used as given; a relative path
  (\code{"100s"}, \code{"data/100s"}), read from the input directory if one is set,
  otherwise from the current directory; or a URL
  (\code{"https://physionet.org/files/mitdb/1.0.0/100"}), read only if network access is
  on (as for \code{read\_record}).
\item[annotators] The annotator names whose files are checked, such as
  \code{\{"atr"\}}. Default \code{\{\}}: none.
\item[input] The input directory for this call only (optional), under which relative
  record names are read (it does not apply to full paths or URLs), and network access
  (\cref{tab:cpp-inputoptions}). Default \code{\{\}}: the current settings.
\end{params}
\returns{A \code{CheckReport} (\cref{tab:cpp-checkreport}): \code{report.ok} is
\code{true} if no problem is an error, \code{report.issues} lists the problems
(\code{Issue} values, \cref{tab:cpp-issue}), and \code{report.to\_string()} gives one line
per problem, or \code{no problems found}.}
\raises{\code{Error} if \code{record} is empty (``record must not be empty''). A record
that cannot be read is reported (code \code{header\_unreadable}), not thrown.}
\begin{lstlisting}[language=C++,style=example]
CheckReport report = check_record("100s", {"atr"});
std::printf("%d\n", report.ok);                    // 1: no errors
std::printf("%s\n", report.to_string().c_str());
// warning: signal 0: the units are not specified
// warning: signal 1: the units are not specified

// A damaged record: one sample of a 12-bit signal is overwritten with 2048.
std::vector<WFDB_Sample> x(100);
for (int i = 0; i < 100; i++) x[i] = i;
WriteOptions w;
w.gains = {100};
w.adcres = {12};
w.units = {"mV"};
w.names = {"ramp"};
write_record("damaged", x, 1, 100, w);
std::FILE *f = std::fopen("damaged.dat", "r+b");
std::fseek(f, 20, SEEK_SET);           // sample 10 (2 bytes per sample in format 16)
std::fputc(0x00, f);                   // 0x0800 = 2048, little-endian
std::fputc(0x08, f);
std::fclose(f);
report = check_record("damaged");
std::printf("%d\n", report.ok);                    // 0
for (const Issue &i : report.issues)
    std::printf("%s %s %s\n", i.severity.c_str(), i.code.c_str(), i.where.c_str());
// error sample_range signal 0
// error checksum signal 0
std::printf("%s\n", report.lines()[0].c_str());
// error: signal 0: 1 sample is outside the ADC range (-2048 to 2047): sample 10 (value 2048)
\end{lstlisting}

% ------------------------------------------------------------------------------------------
\subsection{Header metadata}
\label{sec:cpp-metadata}

These functions read and change the metadata of a header: its info strings of the form
\code{\# Key: value}, such as \code{\# Age: 63} (\cref{chap:metadata}). In short: keys are
compared without regard to letter case (A-Z and a-z only); a key must not be empty and
must not contain \code{:}, \code{\#}, CR, LF or NUL, and a value must not contain CR, LF
or NUL; spaces, tabs, VT, FF and CR at either end of keys and values are removed; keys and
values are bytes (\code{std::string}). A change rewrites the header in one step, through a
temporary file in the directory of the destination, which then replaces it, keeping its
permissions and every other byte of it; all the arguments are checked first, so an error
leaves the header unchanged, and nothing is written if nothing changes (unless
\code{dir} is given). Only the \file{.hea} file is read (not a \file{.info} file); for a
multi-segment record, it is the master header, unless a segment is named. The results
and arguments are of the types \code{Metadata} and \code{MetadataChanges}:
\begin{lstlisting}[language=C++]
using Metadata = std::vector<std::pair<std::string, std::string>>;
using MetadataChanges =
    std::vector<std::pair<std::string, std::optional<std::string>>>;
\end{lstlisting}
Errors throw \code{ArgumentError} (an invalid argument) or \code{FileError} (a file that
cannot be read or written), both subclasses of \code{Error} (\cref{sec:cpp-warnings}).
These functions do not call \code{wfdbquit()}. Metadata often holds information about
patients: WFDB-x stores it as plain text, exactly as given, and does not de-identify it.

\begin{phibox}
A header can hold identifying data about a patient (a name, a date of birth, a medical
record number). The metadata functions read and write such data as given. Records that
are shared must be de-identified first.
\end{phibox}

\apientry{read\_metadata}{Read the metadata of a record}
\begin{lstlisting}[language=C++,style=signature]
Metadata read_metadata(const std::string &record,
                       const InputOptions &input = {});
\end{lstlisting}
Reads the metadata entries of the header of a record, in file order. An info string is an
entry if it contains a colon: the key is the text before the first colon, the value the
rest of the line, both without spaces at either end. Info strings without a colon, or
whose key is empty or contains \code{\#}, are free comments, which are not returned.
\begin{params}
\item[record] Record name, without the \file{.hea} extension (a trailing \code{".hea"} is
  also accepted): a full path, absolute (\code{"/data/mitdb/100"}) or in the home
  directory (\code{"\textasciitilde/Documents/100s"}), used as given; a relative path
  (\code{"100s"}, \code{"data/100s"}), read from the input directory if one is set,
  otherwise from the current directory; or a URL
  (\code{"https://physionet.org/files/mitdb/1.0.0/100"}), read only if network access is
  on.
\item[input] The input directory for this call only (optional), under which relative
  record names are read (it does not apply to full paths or URLs), and network access
  (\cref{tab:cpp-inputoptions}). Default \code{\{\}}: the current settings.
\end{params}
\returns{The entries, as (key, value) pairs, in file order. A key that appears more than
once (ignoring letter case) is returned once, with its first value, with a warning
(\cref{sec:cpp-warnings}).}
\raises{\code{FileError} if the header cannot be found or read (``can't read header for
record REC''); \code{ArgumentError} if \code{record} is empty (``record name must not be
empty'') or \code{"-"} (the standard input).}
\begin{lstlisting}[language=C++,style=example]
WriteOptions w;
w.gains = {200};
w.adcres = {12};
w.units = {"mV"};
w.names = {"ECG"};
w.comments = {"recorded at rest"};
w.metadata = {{"Age", "63"}, {"Sex", "F"}};
write_record("pat", std::vector<WFDB_Sample>(360, 0), 1, 360, w);
for (const auto &e : read_metadata("pat"))
    std::printf("%s = %s\n", e.first.c_str(), e.second.c_str());
// Age = 63
// Sex = F
for (const std::string &c : read_header("pat").comments)
    std::printf("[%s]\n", c.c_str());
// [recorded at rest]: the free comment is not an entry
// [ Age: 63]
// [ Sex: F]
\end{lstlisting}

\apientry{get\_metadata}{Return the value of a metadata entry}
\begin{lstlisting}[language=C++,style=signature]
std::optional<std::string> get_metadata(const std::string &record,
                                        const std::string &key,
                                        const InputOptions &input = {});
\end{lstlisting}
Returns the value of one metadata entry of the header of a record.
\begin{params}
\item[record] Record name (as for \code{read\_metadata}).
\item[key] The key, such as \code{"Age"}. Letter case is ignored, and spaces at either end
  are removed; it must be a valid key (see \code{set\_metadata}).
\item[input] The input directory and network access for this call only. Default
  \code{\{\}}: the current settings.
\end{params}
\returns{The value (\code{""} for an entry with an empty value), or \code{std::nullopt} if
the header has no entry for \code{key}. If it has more than one, the first is returned,
with the warning of \code{read\_metadata}.}
\raises{As \code{read\_metadata}; also \code{ArgumentError} if \code{key} is not valid.}
\begin{lstlisting}[language=C++,style=example]
WriteOptions w;
w.gains = {200};
w.adcres = {12};
w.units = {"mV"};
w.metadata = {{"Age", "63"}, {"Notes", ""}};
write_record("pat", std::vector<WFDB_Sample>(360, 0), 1, 360, w);
std::printf("%s\n", get_metadata("pat", "age")->c_str());       // 63
std::printf("[%s]\n", get_metadata("pat", " NOTES ")->c_str()); // []: an empty value
std::printf("%d\n", get_metadata("pat", "Weight").has_value()); // 0: no such entry
\end{lstlisting}

\apientry{set\_metadata}{Set the value of a metadata entry}
\begin{lstlisting}[language=C++,style=signature]
void set_metadata(const std::string &record, const std::string &key,
                  const std::string &value, const MetadataOptions &opt = {});
\end{lstlisting}
Sets the value of a metadata entry in the header of a record. If the header has an entry
for \code{key} (ignoring letter case), its line is replaced by \code{\# Key: value}, at
the same position and with its line ending, with the spelling of \code{key} given here
(even if the value is the same). Otherwise the entry is added after the last info string
of the header (or, if it has none, after its last line that is not blank), with the line
ending of the record line. Every other byte of the header is kept.
\begin{params}
\item[record] Record name (as for \code{read\_metadata}). The header is changed where it
  was read, unless \code{opt.dir} is given; a header read from a URL can only be written
  elsewhere, with \code{opt.dir}.
\item[key] The key, such as \code{"Age"} or \code{"Date of birth"}. Spaces at either end
  are removed; it must not then be empty, and must not contain \code{:}, \code{\#}, CR,
  LF or NUL.
\item[value] The value, which may be a sentence. Spaces at either end are removed; it
  must not contain CR, LF or NUL, and may be empty (it is then written \code{\# Key:}).
\item[opt] The input directory and network access for reading the header, and the
  directory \code{dir} in which to write the changed header, for this call only
  (\cref{tab:cpp-metadataoptions}). Default \code{\{\}}: the current settings, and the
  header is changed where it was read.
\end{params}
\returns{Nothing.}
\raises{\code{ArgumentError} if \code{key} or \code{value} is not valid; if the header has
more than one entry for \code{key} (``header of record REC has more than one entry for key
KEY (lines L1, L2); remove them first''); if it has no record line; if it was read from a
URL and \code{opt.dir} is not given; or if \code{record} is empty or \code{"-"}.
\code{FileError} if the header cannot be read; if the changed header cannot be written
(``can't write header for record REC''); or if \code{opt.dir} does not exist
(``directory D does not exist'').}
\begin{lstlisting}[language=C++,style=example]
WriteOptions w;
w.gains = {200};
w.adcres = {12};
w.units = {"mV"};
w.metadata = {{"Age", "63"}};
write_record("pat", std::vector<WFDB_Sample>(360, 0), 1, 360, w);
set_metadata("pat", "AGE", "64");                  // replaces the line: "# AGE: 64"
set_metadata("pat", "Date of birth", "1961-02-03");    // adds "# Date of birth: 1961-02-03"
for (const auto &e : read_metadata("pat"))
    std::printf("%s = %s\n", e.first.c_str(), e.second.c_str());
// AGE = 64
// Date of birth = 1961-02-03
try {
    set_metadata("pat", "a:b", "1");
} catch (const ArgumentError &e) {
    std::printf("%s\n", e.what());      // metadata key 'a:b' must not contain ':'
}
\end{lstlisting}

\apientry{remove\_metadata}{Remove a metadata entry}
\begin{lstlisting}[language=C++,style=signature]
int remove_metadata(const std::string &record, const std::string &key,
                    const MetadataOptions &opt = {});
\end{lstlisting}
Removes a metadata entry from the header of a record. Every entry for \code{key}
(ignoring letter case) is removed, with its line ending, so this also repairs a header
that has more than one. Every other byte is kept. Nothing is written if the header has no
entry for \code{key} (unless \code{opt.dir} is given).
\begin{params}
\item[record] Record name (as for \code{set\_metadata}).
\item[key] The key (as for \code{set\_metadata}).
\item[opt] The input directory, network access and output directory \code{dir} of this
  call (as for \code{set\_metadata}). Default \code{\{\}}.
\end{params}
\returns{The number of lines removed (0 if the header has no entry for \code{key}).}
\raises{As \code{set\_metadata}.}
\begin{lstlisting}[language=C++,style=example]
WriteOptions w;
w.gains = {200};
w.adcres = {12};
w.units = {"mV"};
w.metadata = {{"Age", "63"}, {"Sex", "F"}};
write_record("pat", std::vector<WFDB_Sample>(360, 0), 1, 360, w);
std::printf("%d\n", remove_metadata("pat", "sex"));       // 1
std::printf("%d\n", remove_metadata("pat", "Weight"));    // 0: nothing is written
std::printf("%zu\n", read_metadata("pat").size());        // 1
\end{lstlisting}

\apientry{update\_metadata}{Make several changes to the metadata}
\begin{lstlisting}[language=C++,style=signature]
void update_metadata(const std::string &record, const MetadataChanges &changes,
                     const MetadataOptions &opt = {});
\end{lstlisting}
Makes several changes to the metadata of a record in one write. Each (key, value) of
\code{changes} is applied in order, as \code{set\_metadata} does, or, if the value is
\code{std::nullopt}, as \code{remove\_metadata} does. All the changes are checked first:
if one is not valid, none is made.
\begin{params}
\item[record] Record name (as for \code{set\_metadata}).
\item[changes] The changes, as (key, value) pairs; a value of \code{std::nullopt} removes
  the entry. No two keys may differ only in letter case. Empty \code{changes} change
  nothing (and write nothing, unless \code{opt.dir} is given).
\item[opt] The input directory, network access and output directory \code{dir} of this
  call (as for \code{set\_metadata}). Default \code{\{\}}.
\end{params}
\returns{Nothing.}
\raises{As \code{set\_metadata}; also \code{ArgumentError} if two keys differ only in
letter case (``metadata key 'K' is given more than once (keys are compared without regard
to letter case)'').}
\begin{lstlisting}[language=C++,style=example]
WriteOptions w;
w.gains = {200};
w.adcres = {12};
w.units = {"mV"};
w.metadata = {{"Age", "63"}, {"Dx", "164889003"}};
write_record("pat", std::vector<WFDB_Sample>(360, 0), 1, 360, w);
update_metadata("pat", {{"Age", "64"}, {"Dx", std::nullopt}, {"Sex", "F"}});
for (const auto &e : read_metadata("pat"))
    std::printf("%s = %s\n", e.first.c_str(), e.second.c_str());
// Age = 64
// Sex = F

// The header "pat" in the directory "in" is unchanged; the changed copy is written to
// the directory "out" (which must exist).
std::filesystem::create_directories("in");
std::filesystem::create_directories("out");
std::filesystem::copy_file("pat.hea", "in/pat.hea");
MetadataOptions o;
o.input_dir = "in";
o.dir = "out";
update_metadata("pat", {{"Age", "65"}}, o);
std::printf("%s %s\n", get_metadata("in/pat", "Age")->c_str(),
            get_metadata("out/pat", "Age")->c_str());      // 64 65
\end{lstlisting}

% ------------------------------------------------------------------------------------------
\subsection{The options}
\label{sec:cpp-options}

The options of the high-level functions are the fields of the structures of
\cref{tab:cpp-inputoptions,tab:cpp-readoptions,tab:cpp-writeoptions,tab:cpp-metadataoptions}.
\code{ReadOptions} and \code{MetadataOptions} are derived from \code{InputOptions}, and
have its fields too. A field that is not set keeps its default.

\begin{xltabular}{\linewidth}{@{}lL@{}}
\caption{The fields of \code{InputOptions}: where \code{read\_header},
  \code{read\_annotations}, \code{check\_record}, \code{read\_metadata} and
  \code{get\_metadata} find files, for one call}
\label{tab:cpp-inputoptions} \\
\toprule
\tabhead Field & \tabhead Meaning \\
\midrule
\endfirsthead
\toprule
\tabhead Field & \tabhead Meaning \\
\midrule
\endhead
\bottomrule
\endfoot
\code{std::optional<std::string> input\_dir} & the input directory for this call only
  (optional; see \code{wfdbx\_set\_input\_dir}): a local directory, a
  \code{\textasciitilde} path or the URL of a directory, under which relative record names
  are read. It does not apply to full paths or URLs. \code{""} means the default
  (\code{WFDBX\_INPUT\_DIR}, or the current directory). Not set: the current setting \\
\code{std::optional<bool> network} & \code{true} to allow opening URLs for this call only,
  \code{false} to forbid it (see \code{wfdbx\_set\_network}). Not set: the current
  setting \\
\end{xltabular}

\begin{xltabular}{\linewidth}{@{}lL@{}}
\caption{The fields of \code{ReadOptions}, the options of \code{read\_record} (with
  \code{input\_dir} and \code{network} of \code{InputOptions})}
\label{tab:cpp-readoptions} \\
\toprule
\tabhead Field & \tabhead Meaning (default) \\
\midrule
\endfirsthead
\toprule
\tabhead Field & \tabhead Meaning (default) \\
\midrule
\endhead
\bottomrule
\endfoot
\code{std::vector<std::string> channels} & the signals to return, in this order, by number
  (\code{"0"}, \code{"1"}, \ldots) or by name (\code{"V5"}, as in \code{Record::names()},
  then as \code{findsig} finds it). Default empty: all the signals \\
\code{TimeSpec from} & the first sample: a sample number ($\ge 0$) or a time string.
  Default empty: 0 \\
\code{TimeSpec to} & the sample after the last one: a sample number or a time string.
  Default empty: the end of the record \\
\code{bool highres} & \code{true} to read a multi-frequency record at its highest sampling
  frequency, repeating the samples of the slower signals; \code{false} averages the
  samples of the faster signals to one per frame. Default \code{false} \\
\code{bool pad} & \code{true} to replace each invalid sample by the previous valid sample
  of the same signal. Default \code{false} \\
\code{bool checksum\_warning} & \code{true} to print a warning on the standard error when
  the samples do not match the checksums of the header (\cref{sec:cpp-warnings}); the
  samples are returned either way. Default \code{true} \\
\code{bool legacy} & behave as the original WFDB package, for this call only. Default
  \code{false} \\
\end{xltabular}

\begin{xltabular}{\linewidth}{@{}lL@{}}
\caption{The fields of \code{WriteOptions}, the options of \code{write\_record} and
  \code{write\_record\_physical}. Each per-signal vector holds one value for all the
  signals, or one value per signal (or none, for those that have a default)}
\label{tab:cpp-writeoptions} \\
\toprule
\tabhead Field & \tabhead Meaning (default) \\
\midrule
\endfirsthead
\toprule
\tabhead Field & \tabhead Meaning (default) \\
\midrule
\endhead
\bottomrule
\endfoot
\code{int fmt} & the storage format, the same for all the signals, such as \code{16},
  \code{212}, \code{24} or \code{516} (FLAC) (\cref{chap:formats}). Default \code{16} \\
\code{std::vector<WFDB\_Gain> gains} & the gains, in ADC units per physical unit:
  required, and finite. A gain of 0 marks a signal as uncalibrated, which only
  \code{write\_record} (ADC units) allows. Negative gains are allowed \\
\code{std::vector<int> baselines} & the ADC values that correspond to 0 physical units.
  Default: the \code{adczeros} values \\
\code{std::vector<int> adczeros} & the ADC output for an input of 0\,V. Default
  \code{\{0\}} \\
\code{std::vector<int> adcres} & the ADC resolutions, in bits: required, each at least
  1 \\
\code{std::vector<std::optional<std::string>> units} & the physical units, such as
  \code{"mV"}: required for every signal whose gain is not 0; optional for uncalibrated
  signals \\
\code{std::vector<std::string> names} & the signal descriptions. Default: \code{""} \\
\code{std::vector<int> samps\_per\_frame} & the number of samples per frame of each
  signal, for a multi-frequency record (each $\ge 1$). Default: 1 \\
\code{std::optional<std::string> base\_time} & the start time and date, such as
  \code{"12:00:00 01/01/2020"} (see \code{setbasetime}). Default: none \\
\code{std::vector<std::string> comments} & the info strings, written to the header as
  comment lines (\code{"\#"} followed by the string). Default: none \\
\code{Metadata metadata} & metadata entries (\cref{chap:metadata}) written to the header
  after the comments, as \code{\# Key: value} (through \code{putinfo}), in this order. The
  keys and values follow the rules of \code{set\_metadata}; no two keys may differ only
  in letter case, and none may be the key of an entry in \code{comments}
  (\code{ArgumentError}, before any file is created). Default: none \\
\code{std::string dir} & if not empty, the directory (which must exist) in which the
  files are created, for this call only; the record name must then be a plain name,
  without a directory, and the record opens from that directory under its plain name.
  Default \code{""}: the output directory \\
\code{bool legacy} & behave as the original WFDB package, for this call only: the gain
  defaults to 200, and the ADC resolution to one based on the format. Default
  \code{false} \\
\end{xltabular}

\begin{xltabular}{\linewidth}{@{}lL@{}}
\caption{The fields of \code{MetadataOptions}, the options of \code{set\_metadata},
  \code{remove\_metadata} and \code{update\_metadata} (with \code{input\_dir} and
  \code{network} of \code{InputOptions}, which say where the header is read from)}
\label{tab:cpp-metadataoptions} \\
\toprule
\tabhead Field & \tabhead Meaning \\
\midrule
\endfirsthead
\toprule
\tabhead Field & \tabhead Meaning \\
\midrule
\endhead
\bottomrule
\endfoot
\code{std::string dir} & if not empty, the directory (which must exist; a
  \code{\textasciitilde} path is expanded) in which the changed header is written, under
  the file name of the header that was read, instead of changing that header. The header
  is then written even if nothing changes. The signal files are not copied. Default
  \code{""}: the header is changed where it was read \\
\end{xltabular}

\begin{xltabular}{\linewidth}{@{}lL@{}}
\caption{The members of \code{TimeSpec}, a sample number or a time string}
\label{tab:cpp-timespec} \\
\toprule
\tabhead Member & \tabhead Meaning \\
\midrule
\endfirsthead
\toprule
\tabhead Member & \tabhead Meaning \\
\midrule
\endhead
\bottomrule
\endfoot
\code{TimeSpec()} & an empty value: no time is given \\
\code{TimeSpec(WFDB\_Time t)}, \code{TimeSpec(int t)}, \code{TimeSpec(long t)} & the
  sample number \code{t} \\
\code{TimeSpec(const char *s)}, \code{TimeSpec(std::string s)} & the time string
  \code{s} \\
\code{bool empty() const} & \code{true} if no time was given \\
\code{bool needs\_freq() const} & \code{true} if the time is a string that needs a sampling
  frequency to be converted (see \code{wfdbx\_time\_needs\_freq}) \\
\code{WFDB\_Time resolve() const} & the sample number (0 if empty); a time string is
  converted with \code{strtim} for the record open at the time, and a time of day gives
  the positive sample number it corresponds to \\
\end{xltabular}

% ------------------------------------------------------------------------------------------
\subsection{The result types}
\label{sec:cpp-results}

The readers return structures of the types \code{Record}, \code{Header} and
\code{Annotations}, whose members are listed in
\cref{tab:cpp-record,tab:cpp-header,tab:cpp-annotations}; \code{check\_record} returns a
\code{CheckReport}, which holds \code{Issue} values
(\cref{tab:cpp-checkreport,tab:cpp-issue}). The signals of a record are described by
\code{WFDB\_Siginfo} structures (\cref{tab:cpp-siginfo}), which are those of the WFDB
library (\cref{sec:cpp-library}). Their fields are public, and may be changed.

\begin{xltabular}{\linewidth}{@{}lL@{}}
\caption{The members of \code{Record}, the result of \code{read\_record}}
\label{tab:cpp-record} \\
\toprule
\tabhead Member & \tabhead Contents \\
\midrule
\endfirsthead
\toprule
\tabhead Member & \tabhead Contents \\
\midrule
\endhead
\bottomrule
\endfoot
\code{std::string name} & the record name, as given to \code{read\_record} \\
\code{WFDB\_Frequency fs} & the sampling frequency of the returned samples, in Hz, or 0 if
  the header gives none; with \code{highres}, the highest sampling frequency of a
  multi-frequency record \\
\code{int nsig} & the number of signals returned \\
\code{std::vector<int> signals} & the number, in the record, of each signal returned \\
\code{std::vector<WFDB\_Sample> samples} & the samples in ADC units, frame by frame:
  \code{samples[i * nsig + s]} is sample \code{i} of signal \code{s}; invalid samples are
  \code{WFDB\_INVALID\_SAMPLE} \\
\code{std::vector<WFDB\_Siginfo> siginfo} & one per signal returned, in the order of the
  signals; \code{desc} has no white space at either end \\
\code{std::optional<std::string> base\_time} & the start time and date as text, to the
  second, such as \code{"[08:45:00 01/01/2001]"}; \code{std::nullopt} if the header gives
  none \\
\code{std::optional<long long> base\_time\_ms} & the start time in milliseconds after
  midnight (\code{"14:00:03.500"} gives \code{50403500}); given exactly when
  \code{base\_time} is \\
\code{std::optional<std::array<int, 3>> base\_date} & the start date as
  \code{\{year, month, day\}}; \code{std::nullopt} if the header gives no date \\
\code{WFDB\_Time sampfrom} & the number of the first sample returned \\
\code{std::vector<std::string> comments} & the info strings of the record: the comments of
  the header, and the lines of its \file{.info} file if it has one (see \code{getinfo}) \\
\code{bool checksum\_ok} & \code{false} if the samples do not match the checksums of the
  header; checksums are checked only when a read reaches the end of the record, so a
  shorter read gives \code{true} \\
\code{bool legacy} & \code{true} if the record was read in legacy mode; \code{physical()}
  then uses a gain of 200 for uncalibrated signals \\
\code{size\_t nsamples() const} & the number of samples per signal \\
\code{WFDB\_Sample at(size\_t i, int s) const} & sample \code{i} of signal \code{s}, in ADC
  units (not bounds-checked) \\
\code{double physical(size\_t i, int s) const} & sample \code{i} of signal \code{s} in
  physical units, $(\mathit{value} - \mathit{baseline}) / \mathit{gain}$; NaN if the
  sample is invalid. Throws \code{Error} (``signal N of record REC is uncalibrated'') if
  the gain of the signal is 0 \\
\code{std::vector<double> physical\_signal(int s) const} & all the samples of signal
  \code{s} in physical units \\
\code{std::vector<std::string> names() const} & the signal descriptions (\code{""} if
  absent), without the white space that a header may put at either end of them \\
\end{xltabular}

\begin{xltabular}{\linewidth}{@{}lL@{}}
\caption{The members of \code{Header}, the result of \code{read\_header}}
\label{tab:cpp-header} \\
\toprule
\tabhead Member & \tabhead Contents \\
\midrule
\endfirsthead
\toprule
\tabhead Member & \tabhead Contents \\
\midrule
\endhead
\bottomrule
\endfoot
\code{std::string name} & the record name, as given to \code{read\_header} \\
\code{WFDB\_Frequency fs} & the sampling frequency in Hz, as \code{sampfreq()} gives it; 0
  if the header gives none \\
\code{WFDB\_Frequency counter\_freq} & the counter frequency (\code{getcfreq()}) \\
\code{double base\_counter} & the base counter value (\code{getbasecount()}) \\
\code{WFDB\_Time nsamples} & the number of samples per signal, from the record line of the
  header (0 if the header does not give it) \\
\code{int nsig} & the number of signals \\
\code{std::vector<WFDB\_Siginfo> siginfo} & the signal specifications, one per signal;
  \code{desc} has no white space at either end \\
\code{base\_time}, \code{base\_time\_ms}, \code{base\_date} & the start time and date, as
  in \code{Record} \\
\code{std::vector<std::string> comments} & the info strings of the record, as in
  \code{Record} \\
\code{Metadata metadata} & the metadata entries of the header, in file order, as
  \code{read\_metadata} returns them: from the \file{.hea} file only; empty, without an
  error, for an EDF record and for a header read from the standard input \\
\end{xltabular}

\begin{xltabular}{\linewidth}{@{}lL@{}}
\caption{The members of \code{Annotations}, the result of \code{read\_annotations} and the
  argument of \code{write\_annotations}: parallel vectors, annotation \code{i} being
  \code{samples[i]}, \code{codes[i]}, \code{symbols[i]}, \ldots, in file order}
\label{tab:cpp-annotations} \\
\toprule
\tabhead Member & \tabhead Contents \\
\midrule
\endfirsthead
\toprule
\tabhead Member & \tabhead Contents \\
\midrule
\endhead
\bottomrule
\endfoot
\code{std::string record}, \code{annotator} & the record and annotator names, such as
  \code{"100s"} and \code{"atr"} \\
\code{WFDB\_Frequency fs} & the frequency in which \code{samples} are counted, normally the
  sampling frequency of the record; 0 if unknown \\
\code{std::vector<WFDB\_Time> samples} & the annotation times, in samples \\
\code{std::vector<int> codes} & the annotation codes, such as \code{NORMAL} or \code{PVC} \\
\code{std::vector<std::string> symbols} & the mnemonics of the codes, as \code{annstr}
  returns them, such as \code{"N"}, \code{"V"} or \code{"+"} \\
\code{std::vector<int> subtypes}, \code{chans}, \code{nums} & the subtype, channel number
  and annotator number of each annotation \\
\code{std::vector<std::optional<std::string>> aux} & the auxiliary text of each
  annotation, up to its first NUL byte, or \code{std::nullopt} where there is none \\
\code{size\_t size() const} & the number of annotations \\
\end{xltabular}

\begin{table}[htbp]
\centering
\caption{The members of \code{CheckReport}, the result of \code{check\_record}}
\label{tab:cpp-checkreport}
\begin{tabularx}{\linewidth}{@{}lL@{}}
\toprule
\tabhead Member & \tabhead Contents \\
\midrule
\code{std::string record} & the record name, as given \\
\code{bool ok} & \code{true} if no issue has the severity \code{"error"} \\
\code{std::vector<Issue> issues} & the problems, in the order of the checks \\
\code{std::vector<std::string> lines() const} & the text form of the report: one line per
  issue (\code{Issue::to\_string()}), or the one line \code{no problems found} \\
\code{std::string to\_string() const} & the lines of \code{lines()}, joined by line feeds
  (no final line feed) \\
\bottomrule
\end{tabularx}
\end{table}

\begin{table}[htbp]
\centering
\caption{The members of \code{Issue}, a problem found by \code{check\_record}}
\label{tab:cpp-issue}
\begin{tabularx}{\linewidth}{@{}lL@{}}
\toprule
\tabhead Member & \tabhead Contents \\
\midrule
\code{std::string severity} & \code{"error"}, \code{"warning"} or \code{"info"} \\
\code{std::string code} & a stable identifier of the kind of problem, such as
  \code{"checksum"} or \code{"gain\_missing"} (\cref{tab:python-issue-codes}) \\
\code{std::string where} & where the problem is: \code{""} (the record),
  \code{"signal N"}, \code{"segment SEG"}, \code{"segment SEG, signal N"} or
  \code{"annotator ANN"} \\
\code{std::string message} & the description of the problem \\
\code{std::string to\_string() const} & \code{"SEVERITY: WHERE: MESSAGE"}, or
  \code{"SEVERITY: MESSAGE"} when \code{where} is empty \\
\bottomrule
\end{tabularx}
\end{table}

\begin{xltabular}{\linewidth}{@{}lL@{}}
\caption{The fields of \code{WFDB\_Siginfo} (also named \code{Siginfo}), the specification
  of a signal}
\label{tab:cpp-siginfo} \\
\toprule
\tabhead Field & \tabhead Contents \\
\midrule
\endfirsthead
\toprule
\tabhead Field & \tabhead Contents \\
\midrule
\endhead
\bottomrule
\endfoot
\code{std::optional<std::string> fname} & the name of the signal file, as the header gives
  it \\
\code{std::optional<std::string> desc} & the signal description, such as \code{"MLII"} \\
\code{std::optional<std::string> units} & the physical units, such as \code{"mV"};
  \code{std::nullopt} if not specified \\
\code{WFDB\_Gain gain} & the gain, in ADC units per physical unit (0: uncalibrated) \\
\code{WFDB\_Sample initval} & the value of the first sample \\
\code{WFDB\_Group group} & the signal group number (the signals of one group share one
  file) \\
\code{int fmt} & the storage format, such as \code{16}, \code{212} or \code{516} \\
\code{int spf} & the number of samples per frame (more than 1 for the faster signals of a
  multi-frequency record) \\
\code{int bsize} & the block size, in bytes (0 for ordinary files) \\
\code{int adcres} & the ADC resolution, in bits (0: not specified) \\
\code{int adczero} & the ADC output for an input of 0\,V \\
\code{int baseline} & the ADC value that corresponds to 0 physical units \\
\code{long nsamp} & the number of samples (0: not specified) \\
\code{int cksum} & the 16-bit checksum of the samples \\
\end{xltabular}

% ------------------------------------------------------------------------------------------
\subsection{Warnings and errors}
\label{sec:cpp-warnings}

\paragraph{Exceptions.} The high-level functions report errors by throwing exceptions of
the classes of \cref{tab:cpp-exceptions}, all declared in \file{<wfdbx/record.hpp>}.
\code{e.what()} returns the message, which describes the failed operation, such as
``can't read record REC''. Each function checks its arguments before it creates or
changes a file, so an invalid argument leaves the files as they were. The library may
also print a message of its own on the standard error when it cannot open a file (as in
\code{init: can't open header for record REC}); \code{wfdberror()} returns the last such
message, and \code{wfdbquiet()} turns these messages off (\cref{sec:cpp-library}).

\begin{table}[htbp]
\centering
\caption{The exceptions of the high-level API}
\label{tab:cpp-exceptions}
\begin{tabularx}{\linewidth}{@{}lL@{}}
\toprule
\tabhead Exception & \tabhead Thrown when \\
\midrule
\code{Error} (derived from \code{std::runtime\_error}) & a record, header or annotation
  file cannot be found, read, created or written completely (for example on a full
  disk); a signal or annotation file ends before its header or its end-of-file marker
  says; an argument is invalid; or the operation needs a value that the record does not
  give (a time in seconds without a sampling frequency, physical units without a gain) \\
\code{ArgumentError} (derived from \code{Error}) & an argument of a metadata function is
  invalid: an invalid key or value, a record name that is empty or \code{"-"}, a key that
  the header has more than once (\code{set\_metadata}, \code{update\_metadata}), a header
  without a record line, or a header read from a URL to be changed without \code{dir};
  also an invalid \code{WriteOptions::metadata} in \code{write\_record} \\
\code{FileError} (derived from \code{Error}) & a metadata function cannot read or write a
  file: a header that cannot be found or read, a changed header that cannot be written,
  or a \code{dir} that does not exist \\
\bottomrule
\end{tabularx}
\end{table}

\paragraph{Warnings.} Warnings are printed on the standard error with the library's
messages, as \code{wfdb\_error} prints them: \code{wfdbquiet()} silences them, and
\code{wfdberror()} returns the last one. There is no callback.

\paragraph{The checksum warning.} When \code{read\_record} reads to the end of a record
whose samples do not match the checksums of its header, it returns all the samples and
prints the warning ``WFDB-x: checksum error in record REC: the signal data do not match
the checksums in its header''. \code{checksum\_warning = false} turns the warning off;
\code{checksum\_ok} of the result is \code{false} either way.

\paragraph{The duplicate-key warning.} When a header holds two metadata entries whose keys
differ only in letter case (written by other software), \code{read\_metadata},
\code{get\_metadata} and \code{read\_header} use the first and print the warning ``WFDB-x:
header of record REC has more than one entry for key KEY (lines L1, L2); the first is
used''. \code{set\_metadata} and \code{update\_metadata} then refuse to change that key
(\code{ArgumentError}), and \code{remove\_metadata} removes all of its entries.

\begin{lstlisting}[language=C++,style=example]
// A record whose last sample is changed after it was written: the checksum fails.
std::vector<WFDB_Sample> x(100);
for (int i = 0; i < 100; i++) x[i] = i;
WriteOptions w;
w.gains = {100};
w.adcres = {12};
w.units = {"mV"};
write_record("ck", x, 1, 100, w);
std::FILE *f = std::fopen("ck.dat", "r+b");
std::fseek(f, -2, SEEK_END);
std::fputc(0, f);
std::fputc(0, f);
std::fclose(f);
Record rec = read_record("ck");      // prints the checksum warning
std::printf("%zu %d\n", rec.nsamples(), rec.checksum_ok);   // 100 0: all the samples
ReadOptions quiet;
quiet.checksum_warning = false;
rec = read_record("ck", quiet);      // no warning

// Two entries for one key: the first is used, with a warning.
WriteOptions d = w;
d.gains = {0};
d.units = {};
d.comments = {" Age: 63", " AGE: 64"};
write_record("dup", std::vector<WFDB_Sample>(10, 0), 1, 100, d);
std::printf("%s\n", get_metadata("dup", "age")->c_str());   // 63, with a warning
try {
    set_metadata("dup", "Age", "65");
} catch (const ArgumentError &e) {
    std::printf("ArgumentError: %s\n", e.what());
    // ArgumentError: header of record dup has more than one entry for key Age
    // (lines 3, 4); remove them first
}
std::printf("%d\n", remove_metadata("dup", "age"));          // 2

try {
    read_record("nosuchrecord");
} catch (const Error &e) {
    std::printf("%s\n", e.what());   // can't read record nosuchrecord
}
try {
    read_metadata("nosuchrecord");
} catch (const FileError &e) {
    std::printf("%s\n", e.what());   // can't read header for record nosuchrecord
}
\end{lstlisting}

% ==========================================================================================
\section{The WFDB library functions in C++}
\label{sec:cpp-library}

Every function of the WFDB library is declared in \file{<wfdbx/wfdb.hpp>}, in namespace
\code{wfdbx}, under its C name, with the same arguments, in the same order, and the same
return values, so that programs and documentation written for the C library carry over
directly. The functions, their arguments and their results are documented once, for all
the languages, in \cref{chap:library}; this section gives the rules by which the C
declarations become C++ functions (see also \cref{sec:library-cpp}). The header
summarizes each function, and marks what WFDB-x changes in it.

\subsection{Conventions}

\paragraph{Names and headers.} The functions, the types (\code{WFDB\_Sample},
\code{WFDB\_Time}, \ldots), the structures and the constants (\code{WFDB\_READ},
\code{WFDB\_INVALID\_SAMPLE}, \code{WFDB\_HIGHRES}, \ldots) keep their C names, in
namespace \code{wfdbx}. \file{<wfdbx/ecgcodes.hpp>} defines the annotation codes
(\code{NORMAL}, \code{PVC}, \ldots, \cref{chap:formats}), and \file{<wfdbx/ecgmap.hpp>}
the macros of \file{<wfdb/ecgmap.h>} as inline functions (\code{isann}, \code{isqrs},
\code{map1}, \ldots, which call \code{wfdb\_isann}, \code{wfdb\_isqrs},
\code{wfdb\_map1}, \ldots). \file{<wfdbx/record.hpp>} includes \file{<wfdbx/wfdb.hpp>}.

\paragraph{Arguments.} String arguments are \code{const char *}, and \code{nullptr} is
accepted where C accepts \code{NULL}: \code{sampfreq(nullptr)} returns the frequency of
the open record, and \code{getinfo(nullptr)} continues with the info strings of the last
record. Arrays and structures are allocated by the caller and passed by pointer, as in
C: \code{isigopen} fills an array of \code{WFDB\_Siginfo}, \code{getvec} an array of
\code{WFDB\_Sample}, \code{getann} a \code{WFDB\_Annotation}. \code{getseginfo} takes a
\code{const WFDB\_Seginfo **}.

\paragraph{Strings} are returned by value, as \code{std::string}, or as
\code{std::optional<std::string>} where the C function may return \code{NULL}
(\code{annstr}, \code{ecgstr}, \code{anndesc}, \code{getinfo}, \code{wfdbfile}, \ldots);
unlike C's pointers to static buffers, they stay valid after later calls. Text is handled
as bytes.

\paragraph{Return values and errors.} The functions return what the C functions return:
an \code{int} status (negative on errors), a number of signals, a \code{double}, a string
or \code{std::nullopt}. Errors are reported by these values and by a message on the
standard error, as in C, never by exceptions. \code{wfdberror()} returns the last message;
\code{wfdbquiet()} turns the messages off and \code{wfdbverbose()} on again.

\paragraph{State.} As in C, the library keeps its state in the process: the open record,
its input signals, the output signals and the annotators remain open between calls, until
\code{wfdbquit()} closes them. \code{wfdbquit()} leaves the input and output directories
and the network setting as they are. The library is not thread-safe. The high-level
functions of \cref{sec:cpp-api} (except the metadata functions) call \code{wfdbquit()},
which closes anything opened with the library functions.

\paragraph{Additions.} \code{getvecs(buffer, n)} reads up to \code{n} sample vectors at
once, with exactly the result of \code{n} calls of \code{getvec}, but decoding whole
blocks of the signal files; \code{getvecs(buffer, n, \&status)} also gives the status of
the last \code{getvec} call:
\begin{lstlisting}[language=C++]
long long getvecs(WFDB_Sample *out, long long n, int *status = nullptr);
\end{lstlisting}

\subsection{The structures}
\label{sec:cpp-structures}

The structures of the C library have the names and the fields of the C structures,
with C's zero values as defaults; they are also available under the shorter names
\code{Siginfo}, \code{Anninfo}, \code{Annotation}, \code{Calinfo} and \code{Seginfo}
(\cref{tab:cpp-structures}).

\begin{xltabular}{\linewidth}{@{}lL@{}}
\caption{The C structures in C++}
\label{tab:cpp-structures} \\
\toprule
\tabhead Structure & \tabhead Fields \\
\midrule
\endfirsthead
\toprule
\tabhead Structure & \tabhead Fields \\
\midrule
\endhead
\bottomrule
\endfoot
\code{WFDB\_Siginfo} & \code{fname}, \code{desc}, \code{units} (\code{std::optional<std::string>}),
  \code{gain}, \code{initval}, \code{group}, \code{fmt}, \code{spf}, \code{bsize},
  \code{adcres}, \code{adczero}, \code{baseline}, \code{nsamp}, \code{cksum}
  (\cref{tab:cpp-siginfo}) \\
\code{WFDB\_Seginfo} & \code{recname} (the name of the segment, a \code{char} array),
  \code{nsamp} (its number of samples), \code{samp0} (the number of its first sample in
  the record) \\
\code{WFDB\_Anninfo} & \code{name} (the annotator name, a \code{std::string}), \code{stat}
  (\code{WFDB\_READ} or \code{WFDB\_WRITE} for the MIT format, \code{WFDB\_AHA\_READ} or
  \code{WFDB\_AHA\_WRITE} for the AHA format) \\
\code{WFDB\_Annotation} & \code{time} (\code{WFDB\_Time}, in sample intervals),
  \code{anntyp} (\code{char}), \code{subtyp} (\code{signed char}), \code{chan}
  (\code{unsigned char}), \code{num} (\code{signed char}), \code{aux}
  (\code{std::optional<std::string>}: absent, or a string whose first byte is the length
  of the text that follows, as in C) \\
\code{WFDB\_Calinfo} & \code{low}, \code{high} (the levels of the calibration pulse, in
  physical units), \code{scale} (physical units per cm), \code{sigtype} and \code{units}
  (\code{std::optional<std::string>}), \code{caltype} \\
\end{xltabular}

\code{WFDB\_Time} is \code{long long} (64 bits), \code{WFDB\_Sample} is \code{int},
\code{WFDB\_Date} is \code{long}, \code{WFDB\_Frequency} and \code{WFDB\_Gain} are
\code{double}, and \code{WFDB\_Group}, \code{WFDB\_Signal} and \code{WFDB\_Annotator} are
\code{unsigned int}.

\subsection{The settings functions}
\label{sec:cpp-settings}

The functions of \cref{tab:cpp-settings} are additions of WFDB-x to the library, declared
in \file{<wfdbx/wfdb.hpp>}: they set where files are read and written, network access and
legacy mode, for the high-level functions, the library functions and the tools alike
(\cref{chap:files,chap:behavior}). Each \code{set} function returns the previous setting,
so that it can be restored. The C function \code{setwfdb(p)} sets the input directory too
(\code{p} must name at most one directory), and \code{getwfdb()} returns it (\code{"."} if
none is set); in legacy mode they set and return the WFDB path, as in C.

\begin{xltabular}{\linewidth}{@{}lL@{}}
\caption{The settings functions}
\label{tab:cpp-settings} \\
\toprule
\tabhead Function & \tabhead Effect and result \\
\midrule
\endfirsthead
\toprule
\tabhead Function & \tabhead Effect and result \\
\midrule
\endhead
\bottomrule
\endfoot
\code{std::optional<std::string> wfdbx\_set\_input\_dir(const char *dir)} & sets the input
  directory: a local directory, a \code{\textasciitilde} path or a URL; \code{nullptr} or
  \code{""} restores the default (\code{WFDBX\_INPUT\_DIR}, or else the current
  directory). Returns the previous input directory, or \code{std::nullopt} if none was
  set \\
\code{std::optional<std::string> wfdbx\_get\_input\_dir()} & returns the input directory,
  or \code{std::nullopt} if none is set \\
\code{std::optional<std::string> wfdbx\_set\_output\_dir(const char *dir)} & sets the
  output directory: a local directory or a \code{\textasciitilde} path (a URL is refused,
  with a message); \code{nullptr} or \code{""} restores the default
  (\code{WFDBX\_OUTPUT\_DIR}, or else the current directory). Returns the previous output
  directory, or \code{std::nullopt} \\
\code{std::optional<std::string> wfdbx\_get\_output\_dir()} & returns the output directory,
  or \code{std::nullopt} if none is set \\
\code{int wfdbx\_set\_network(int on)} & turns network access on (1) or off (0); a
  negative number restores the default (on if \code{WFDBX\_NETWORK} is a nonzero number).
  Returns the previous setting, 1 or 0 \\
\code{int wfdbx\_get\_network()} & returns 1 if network access is on, 0 otherwise \\
\code{int wfdbx\_set\_legacy(int on)} & turns legacy mode on (1) or off (0), for the whole
  library. Returns the previous setting \\
\code{int wfdbx\_get\_legacy()} & returns 1 if legacy mode is on, 0 otherwise \\
\code{int wfdbx\_legacy\_args(int \&argc, char **\&argv)} & if \code{argv[1]} is
  \code{"-{}-legacy"}, removes it from the arguments, turns legacy mode on and returns 1;
  otherwise returns 0. It also installs the exit check of the tools (an exit status of 0
  becomes 1 after a write error, outside legacy mode) \\
\code{int wfdbx\_time\_needs\_freq(const char *time\_string)} & returns 1 if
  \code{strtim(time\_string)} needs a sampling frequency, 0 otherwise \\
\code{int wfdbx\_write\_errors()} & returns the number of write errors since the program
  started (output files that could not be written, flushed or closed) \\
\code{void wfdbx\_note\_write\_error()} & counts a write error, for a program that writes
  a file itself \\
\code{std::optional<std::string> wfdbx\_input\_path(const char *name)} & returns the path
  under which a program opens a file that it reads itself, where \code{wfdb\_open} finds a
  file of that name: in normal mode, a relative name under the input directory, a
  \code{\textasciitilde} name expanded, an absolute name as given; in legacy mode, the name
  as given. \code{""} and \code{"-"} are returned as given; \code{std::nullopt} for
  \code{nullptr} or a URL (also once the input directory is prefixed) \\
\code{std::optional<std::string> wfdbx\_output\_path(const char *name)} & returns the path
  under which a program creates a file that it writes itself, where \code{wfdb\_open}
  creates a file of that name: in normal mode, a relative name under the output directory,
  a \code{\textasciitilde} name expanded, an absolute name as given; in legacy mode, any
  name under the output directory, as a plain prefix. The missing directories are created
  when the output directory is prefixed. \code{""} and \code{"-"} are returned as given;
  \code{std::nullopt} for \code{nullptr} or a URL \\
\end{xltabular}

\subsection{Examples}

Reading the samples and the annotations of a record, as a C program does (the record
\code{100s} is read from the current directory):

\begin{lstlisting}[language=C++,style=example]
#include <cstdio>
#include <vector>
#include <wfdbx/ecgcodes.hpp>
#include <wfdbx/wfdb.hpp>
using namespace wfdbx;

int main() {
    WFDB_Siginfo si[2];                       // receives one WFDB_Siginfo per signal
    int nsig = isigopen("100s", si, 2);       // open 2 signals
    std::printf("%d %s %d %g\n", nsig, si[0].desc->c_str(), si[0].fmt, si[0].gain);
    // 2 MLII 212 200
    std::printf("%g\n", sampfreq(nullptr));   // 360: the frequency of the open record

    WFDB_Sample v[2];                         // receives one sample per signal
    isigsettime(strtim("0:30"));              // skip to 30 seconds
    getvec(v);
    std::printf("%d %d %g\n", v[0], v[1], aduphys(0, v[0]));    // 947 968 -0.385

    std::vector<WFDB_Sample> frames(360 * 2); // the next second, in one call
    int status;
    long long n = getvecs(frames.data(), 360, &status);
    std::printf("%lld %d\n", n, status);      // 360 2

    WFDB_Anninfo ai{"atr", WFDB_READ};
    WFDB_Annotation a;
    annopen("100s", &ai, 1);
    int count = 0;
    while (getann(0, &a) == 0) {              // 0: an annotation was read; -1: the end
        if (count < 2)
            std::printf("%s %s\n", mstimstr(a.time).c_str(), annstr(a.anntyp)->c_str());
        count++;
    }
    std::printf("%d annotations\n", count);   // 75 annotations

    wfdbquiet();                              // no messages on the standard error
    int stat = isigopen("nosuchrecord", nullptr, 0);
    std::printf("%d %s", stat, wfdberror().c_str());
    // -1 init: can't open header for record nosuchrecord
    wfdbverbose();
    wfdbquit();                               // close everything
}
\end{lstlisting}

Writing a record and an annotation file, sample by sample and annotation by annotation:

\begin{lstlisting}[language=C++,style=example]
#include <cstdio>
#include <wfdbx/ecgcodes.hpp>
#include <wfdbx/wfdb.hpp>
using namespace wfdbx;

int main() {
    WFDB_Siginfo out;
    out.fname = "ramp.dat";
    out.desc = "ramp";
    out.units = "mV";
    out.gain = 200;
    out.fmt = 16;
    out.adcres = 12;
    std::printf("%d\n", osigfopen(&out, 1));    // 1: one output signal
    setsampfreq(250);
    for (WFDB_Sample k = 0; k < 500; k++) putvec(&k);
    std::printf("%d\n", newheader("ramp"));     // 0: ramp.hea is written
    WFDB_Anninfo ai{"qrs", WFDB_WRITE};
    annopen("ramp", &ai, 1);
    for (WFDB_Time t : {100, 350}) {
        WFDB_Annotation a;
        a.time = t;
        a.anntyp = NORMAL;
        putann(0, &a);
    }
    wfdbquit();                                 // completes ramp.dat and ramp.qrs
}
\end{lstlisting}

% ==========================================================================================
\section{The command-line tools in C++}
\label{sec:cpp-tools}

The 33 command-line tools (\code{rdsamp}, \code{rdann}, \code{wrsamp}, \code{xform},
\ldots) are separate programs, built in \file{cpp/build/bin/} (or the \file{bin} directory
of the installation), one per file of \file{cpp/apps}, and run by their names. Each
follows its C program line by line. The tools, their options, their output and their exit
status are those of the C tools, and are documented in \cref{chap:tools}.

\begin{lstlisting}[style=shell]
export PATH=$PWD/cpp/build/bin:$PATH
export WFDBX_INPUT_DIR=tests/data     # run from the root of the repository
rdsamp -r 100s -t 10 -p     # the first 10 seconds of 100s, in physical units
rdann -r 100s -a atr        # its annotations
rdsamp --legacy -r 100s     # legacy mode: --legacy as the first argument
\end{lstlisting}

There is no launcher: a tool is run by its own name, and it has no \opt{-{}-version}
option, like the C tools. When an output file or the standard output cannot be written (a
full disk, for example), a tool reports it (``PROGRAM: write error on standard output'')
and exits with status 1; with \opt{-{}-legacy}, it keeps the original exit status.

\paragraph{Environment variables.} The tools, like the rest of the library, use the
environment variables of \cref{tab:cpp-environment}, also for the files that they read or
create themselves: the \opt{-i} file of \code{wrsamp}, \code{edf2mit}, \code{wav2mit},
\code{a2m} and \code{ad2m}, the input files of \code{ahaecg2mit} and the \opt{-S} script
of \code{xform} are read from the input directory, and the WAV file of \code{mit2wav},
the EDF file of \code{mit2edf}, the master header of \code{snip -m} and the \file{.des},
\file{.ecg}, \file{.key} and \file{.txt} files of \code{parsescp} are created in the
output directory, as the records they read and write are (with
\code{wfdbx\_input\_path} and \code{wfdbx\_output\_path}, \cref{tab:cpp-settings}). Such
a file cannot be read from a URL.

\begin{table}[htbp]
\centering
\caption{The environment variables used by the C++ library}
\label{tab:cpp-environment}
\begin{tabularx}{\linewidth}{@{}lL@{}}
\toprule
\tabhead Variable & \tabhead Meaning \\
\midrule
\code{WFDBX\_INPUT\_DIR} & the input directory (a directory, a \code{\textasciitilde}
  path or a URL), when \code{wfdbx\_set\_input\_dir} does not set one \\
\code{WFDBX\_OUTPUT\_DIR} & the output directory, when \code{wfdbx\_set\_output\_dir} does
  not set one \\
\code{WFDBX\_NETWORK} & network access is on if it is a nonzero number, when
  \code{wfdbx\_set\_network} does not set it \\
\code{WFDBANNSORT} & \code{0} to keep annotation files in the order written (no sorting
  when they are closed) \\
\code{WFDBGVMODE} & the default getvec mode (\code{setgvmode} with a negative mode) \\
\code{WFDBCAL} & the calibration file read by \code{calopen(nullptr)} (default
  \file{wfdbcal}) \\
\code{WFDB} & the WFDB path, used in legacy mode only (outside legacy mode, it is
  ignored, with a message) \\
\code{HOME} (\code{USERPROFILE} on Windows) & the home directory, for
  \code{\textasciitilde} paths \\
\bottomrule
\end{tabularx}
\end{table}

Tools that write signal parameters require them as options: \code{wrsamp} (with
\opt{-o}) requires \opt{-F} (frequency), \opt{-G} (gain), \opt{-R} (ADC resolution) and,
for calibrated signals, \opt{-U} (units); \code{wav2mit} requires \opt{-G}, and \opt{-U}
unless the gain is 0; \code{ad2m} and \code{ahaecg2mit} require \opt{-F}, \opt{-G},
\opt{-R}, and \opt{-U} unless the gain is 0; \code{a2m}, \code{rdann}, \code{ann2rr},
\code{sortann} and \code{m2a} require \opt{-F} when the record's header gives no
frequency (\cref{chap:tools}).

% ==========================================================================================
\section{Testing}
\label{sec:cpp-testing}

The unit tests of the C++ implementation are in \file{cpp/tests}, built with the library
(unless \code{WFDBX\_BUILD\_TESTS} is \code{OFF}) as the program
\file{cpp/build/wfdbx\_tests}. From the root of the repository, after building:

\begin{lstlisting}[style=shell]
ctest --test-dir cpp/build                          # the unit tests
cpp/build/wfdbx_tests                               # the same, one line per test
cpp/build/wfdbx_tests metadata_shared_cases         # selected tests, by name
\end{lstlisting}

The unit tests cover the WFDB library functions, the high-level functions, the metadata
functions (with the cases shared by the seven languages, in \file{tests/metadata}),
\code{check\_record} and the FLAC codec. They run in temporary copies of the test records
of \file{tests/data} (the environment variable \code{WFDBX\_TEST\_DATA} names another
directory). The tests of network access use a local HTTP server on 127.0.0.1 only, and
are skipped without \cmd{curl}. The tests of \code{check\_record} write the corpus of
\file{python/tests/check\_corpus.py} with \cmd{python3} (or the program named by the
environment variable \code{WFDBX\_PYTHON}) and the Python package of \file{python/src},
and compare the reports with its \file{expected.txt}; without Python they are skipped.
\cmd{cpp/build/wfdbx\_check\_corpus DIR} prints the reports for a corpus in \file{DIR}.

The cross-language conformance suite runs every tool and compares its results, byte for
byte, with those of the original C programs (1135 cases), and, with \opt{-{}-legacy}, in
legacy mode (883 cases). It runs the tools of \file{cpp/build/bin}:

\begin{lstlisting}[style=shell]
PYTHONPATH=$PWD/python/src python3 tests/conformance/harness.py check --impl cpp -j8
PYTHONPATH=$PWD/python/src python3 tests/conformance/harness.py check --impl cpp -j8 --legacy
\end{lstlisting}
